\chapter{Introducción}

\section{Contexto y problemática}
Una tienda de electrónica pequeña administraba su operación con planillas y cuadernos: el conteo del inventario se hacía a mano, las ventas no dejaban constancia utilizable y no existía forma fiable de consultar una garantía o saber quién abrió la caja. Los costos de ese esquema son conocidos: descuadres de stock, ventas sin respaldo, garantías imposibles de rastrear y cero trazabilidad de la jornada.

\section{¿Qué es Tech-Flow?}
Tech-Flow es una aplicación de escritorio para la gestión operativa de una tienda de electrónica. Está escrita en Java 17 con JavaFX 17.0.2, se ejecuta con Maven y persiste sus datos en archivos CSV legibles por cualquier hoja de cálculo. La interfaz organiza las funciones por rol: administrador, cajero y almacenero. Una sola ventana aloja todas las pantallas; los avisos, formularios y confirmaciones aparecen como cuadros dentro de la misma ventana, lo que la hace amistosa para gestores de ventanas en mosaico (p. ej. bspwm).

\section{Propósito y contribución}
Automatizar el núcleo de la operación de la tienda con un sistema verificable: autenticación por rol con selección explícita, caja que bloquea o habilita las ventas, inventario con alertas de stock crítico, ventas con garantías asociadas, entradas de almacén con vista previa y exportación de informes de texto. El proyecto además sirve como evidencia técnica académica: suite de pruebas, documentación y decisiones de diseño registradas.

\section{Alcance y limitaciones}
El sistema cubre autenticación, caja, catálogo e inventario, entradas y pedidos, ventas y proformas, garantías, clientes y proveedores. No incluye base de datos relacional, conexión en red, múltiples sucursales ni respaldo automático: la persistencia es en archivos CSV en la carpeta \texttt{data/} del proyecto. El archivo de empleados guarda contraseñas sin cifrar, lo que se documenta como limitación conocida en el informe de cambios del repositorio.
